\emph{Setting and imports.} Work in $\Cl(1,3)$ with $\gamma_{0}^{2} = 1$, $\gamma_{i}^{2} = -1$ and $I_{4} = \gamma_{0}\gamma_{1}\gamma_{2}\gamma_{3}$. A rotor $R$ ($R\tilde R = 1$, even) acts on the series' Dirac values by $\psi \mapsto R\psi$ (Proposition \claimref{field-theory.second-quantisation}), on vector values (potentials, gradients) and bivector values by $x \mapsto Rx\tilde R$, and trivially on scalar values (the components of any number of scalar doublets and the ghosts of \claimref{synthesis.self-adjoint-brst-charge-needs-a-krein-form-and-the-tower-supplies-the-ghosts}). The covering group of the Lorentz group is the rotor group; its element $-1$ is the end point of the lift of a closed loop of rotations. The free fields have the statistics of \claimref{synthesis.locality-of-the-attached-free-fields-fixes-their-statistics}: with positive norm the spinor is a CAR field and the scalars and the field strength CCR fields; the ghosts are local CAR scalars, which needs the indefinite form. The potential itself in Feynman gauge is not covered by that result; it enters as a condition (Step 2). From outside the store the proof uses only the axioms and results of causal perturbation theory (standard): time-ordered products $T$ that are multilinear, graded symmetric in their arguments and preserve the Grassmann parity (the renormalisation freedom is spanned by local terms of the parity of the product), satisfy the Wick expansion in the free fields ((N3) of \cite{DutschFredenhagen1999}), factorise causally ($T(A_{1}(x_{1})\cdots A_{n}(x_{n})) = T(A_{1}\cdots A_{k})T(A_{k+1}\cdots A_{n})$ when no $x_{j}$, $j > k$, lies in the closed causal future of an $x_{i}$, $i \le k$) and are covariant, $\mathrm{Ad}\,U(L)\,T(A_{1}(x_{1})\cdots) = T(\mathrm{Ad}\,U(L)A_{1}(x_{1})\cdots)$ \cite{DutschFredenhagen1999} ((2.2), (2.3)); there (N1) is stated for the proper orthochronous Poincar\'e group, and covariance under its covering group, with $U$ acting on spinor fields through the rotor group, is an assumption of the present result (standard for spinor fields, cited from memory); their existence \cite{EpsteinGlaser1973} (Fact \claimref{synthesis.imported-epstein-glaser-extension-theorem-and-ward-preserving-time-ordered-products}(i)); Bogoliubov's interacting fields \cite{DutschFredenhagen1999} ((2.7)); and graded locality of interacting fields to all orders \cite{DutschFredenhagen1999} (Proposition 2 and (2.25), read: the bracket is an anticommutator only when both fields have an odd number of spinor fields or both an odd number of ghost fields, in a convention where free spinors commute with free ghosts; in their QED the ghosts do not couple); with a coupled ghost vertex (Yang-Mills) \cite{Hollands2008} (cited from memory, the relevant sections not read).

\emph{Step 1: the rotation by $2\pi$ (clause (i), first sentence).} Let $B = \gamma_{1}\gamma_{2}$, $B^{2} = -1$. The versor $q = 1 + B$ has $q\tilde q = 2$ and $q x \tilde q/2$ is the quarter turn in the plane $B$; $q^{2} = 2B$, so $q^{2}/2 = B$ is the half turn ($B\gamma_{1}\tilde B = -\gamma_{1}$, checked), and $q^{4} = (2B)^{2} = -4$. The four quarter turns compose the closed loop of rotations by $\pi/2, \pi, 3\pi/2, 2\pi$; its lift to rotors, normalised by $(q\tilde q)^{2} = 4$, ends at $q^{4}/4 = -1$. The script computes $q^{4}$, its norm and the action of $U = -1$: $Ux\tilde U = x$ on five vectors and two bivectors, $U\psi = -\psi$ on three even multivectors. So $U(2\pi)$ multiplies a monomial by $(-1)^{n_{\psi}}$, $n_{\psi}$ its number of spinor factors (the conjugate $\tilde\psi$ counts as a spinor factor, $\tilde\psi \mapsto \tilde\psi\tilde R$).

\emph{Step 2: the interaction density (clause (i)).} The monomials are those of the stored results: the vector vertex $\langle J A\rangle_{0}$, $J = \psi\gamma_{0}\tilde\psi$, of the one-loop QED result (\claimref{synthesis.one-loop-qed-in-the-spacetime-algebra}, the matrix reading $\bar\psi\gamma^{\mu}\psi A_{\mu}$); the axial vertex $\langle K A\rangle_{0}$, $K = \psi\gamma_{3}\tilde\psi$ (with the vector vertex it gives the chiral couplings of the weak factor); the Yukawa terms of the doublets and their mass insertions $\langle\psi\tilde\psi\rangle_{0}\phi$, $\langle\psi\tilde\psi I_{4}\rangle_{0}\phi$; the scalar gauge couplings $\langle A\nabla\rangle_{0}\phi\phi$ and $\langle AA\rangle_{0}\phi\phi$ and the quartic $\phi^{4}$; the non-abelian cubic $\langle(\nabla \wedge A)(A \wedge A)\rangle_{0}$ and quartic $\langle(A \wedge A)(A \wedge A)\rangle_{0}$ with structure constants as scalar coefficients; the gauge fixing $\langle\nabla A\rangle_{0}^{2}$; and the ghost terms $\langle\nabla\bar c\,\nabla c\rangle_{0}$ and $\langle\nabla\bar c\,A\rangle_{0}c$ of the BRST result. The scalar content enters only through Lorentz scalars of scalar values, so the number of doublets plays no role. Each is a scalar part of a product of rotor-covariant values, hence invariant under every rotor: $\psi\gamma_{\mu}\tilde\psi \mapsto R\psi\gamma_{\mu}\tilde\psi\tilde R$, $\psi\tilde\psi \mapsto R\psi\tilde\psi\tilde R = \psi\tilde\psi$ (its scalar and $I_{4}$ parts are invariant), and $\langle R X\tilde R\rangle_{0} = \langle X\rangle_{0}$. The script evaluates every monomial on three spinors and sample vectors and scalars, before and after four exact rotors (a rotation $(3 + 4\gamma_{1}\gamma_{2})/5$, a boost $(5 + 4\gamma_{0}\gamma_{1})/3$, a rotation $(5 + 12\gamma_{2}\gamma_{3})/13$ and their product), and under $U = -1$, where the value is $(-1)^{n_{\psi}}$ times the original and nonzero; it also checks that a one-spinor scalar such as $\langle\psi\gamma_{0}A\rangle_{0}$ changes sign under $U$, so the test is not vacuous. The counts are not labelled by hand: the script computes $n_{\psi}$, $n_{c}$, $n_{\bar c}$ as the homogeneity degrees of each monomial under $\psi \mapsto 2\psi$, $c \mapsto 3c$, $\bar c \mapsto 5\bar c$, and finds $n_{\psi} \in \{0, 2\}$, $n_{c} = n_{\bar c}$; the control scalar has degree $1$. The free potential in Feynman gauge: by the Gupta-Bleuler form of the BRST result, $[a_{\mu}, a_{\nu}^{\dagger}] = -\gamma_{\mu}\cdot\gamma_{\nu}$, so its polarisation sum on test vectors is $\sum_{\mu}(-Q(\gamma_{\mu}))(a\cdot\gamma_{\mu})(\gamma_{\mu}\cdot b) = -a\cdot b$ (checked on five vectors), independent of $k$, of degree $d = 0$, and indefinite (negative on a timelike test vector). With the bracket factor $1 + \eta c(-1)^{d}$ of the stored locality result and $c = 1$ (the antiparticle sum equals the particle sum for a Hermitian field), the factor vanishes for $\eta = -1$: the potential is a local commuting field in the indefinite form, the standard Gupta-Bleuler statement (the stored result treats only the field strength at positive norm, so this is a condition of the present result, clause (K1) of Premise \claimref{synthesis.krein-form-premise-indefinite-inner-product-for-gauge-fields}, with the degree computed).

\emph{Step 3: Wick monomials (clause (ii)).} The Grassmann parity of a Wick monomial is the parity of its number of anticommuting factors, $n_{\psi} + n_{c} + n_{\bar c}$, by the free statistics; $U(2\pi)$ acts as $(-1)^{n_{\psi}}$ by Step 1. They agree if and only if $n_{c} + n_{\bar c}$ is even; at ghost number $n_{c} - n_{\bar c} = 0$ the sum $2n_{c}$ is even. The script models the anticommuting factors as tower modes $v_{j} = \tfrac12(e_{j} + f_{j})$ on distinct tower planes of $\Cl(m, m)$ ($e_{j}^{2} = 1$, $f_{j}^{2} = -1$; $v_{j}^{2} = 0$, $\{v_{j}, v_{k}\} = 0$, the CAR of Fact \claimref{field-theory.imported-tower-limit}), reads the Grassmann parity from the grade involution there, and computes the $2\pi$ eigenvalue separately in $\Cl(1,3)$ as the sign by which $U$ applied to each factor changes the product $\psi_{1}\gamma_{0}\psi_{2}\gamma_{0}\cdots\psi_{n_{\psi}}$ of sample spinors (nonzero); over $n_{\psi} \le 3$, $n_{c}, n_{\bar c} \le 2$ the two disagree on exactly the $16$ monomials with $n_{c} + n_{\bar c}$ odd. For ghost-even monomials on disjoint modes the script checks that they anticommute exactly when both have $U(2\pi) = -1$ and commute otherwise. In the convention of \cite{DutschFredenhagen1999}, where spinors commute with ghosts, the bracket sign is $(-1)^{n_{\psi}(A)n_{\psi}(B) + n_{gh}(A)n_{gh}(B)}$; for fields of even ghost count it is $(-1)^{n_{\psi}(A)n_{\psi}(B)}$ in both conventions.

\emph{Step 4: interacting fields keep parity and spin (clause (iii), first part).} The $n$-th order of $A_{\mathrm{int}}(x)$ is $\tfrac{I_{4}^{n}}{n!}\int R(\mathcal{L}(y_{1})\cdots\mathcal{L}(y_{n}); A(x))\,g(y_{1})\cdots g(y_{n})$, a retarded product, a finite combination of time-ordered and anti-time-ordered products of the same factors. By the causal Wick expansion its terms are Wick monomials obtained from the product $\mathcal{L}^{n}A$ by contracting pairs of a field with its conjugate (a $\psi$ with a $\tilde\psi$, a $c$ with a $\bar c$, two bosons), which changes $n_{\psi}$ and $n_{c} + n_{\bar c}$ by even numbers; local renormalisation terms have the Grassmann parity of the product by the parity condition. Since each copy of $\mathcal{L}$ is even (Step 2), every term has the Grassmann parity of $A$. (No conservation of ghost number by the time-ordered products is needed: the result uses only the Grassmann parity and the $2\pi$ eigenvalue, whose ratio is the ghost parity.) By covariance, $\mathrm{Ad}\,U(2\pi)$ acts on the retarded product as on the product of its arguments, that is by $(+1)^{n}(-1)^{n_{\psi}(A)}$. So $A_{\mathrm{int}}$ has the Grassmann parity and the $U(2\pi)$ eigenvalue of $A$, and by Step 3 they agree when $A$ has an even ghost count, in particular when $A$ has ghost number zero. The script checks in the tower model that for an even element $\mathcal{L}$ (spinor and ghost bilinears) the graded commutator $[\mathcal{L}, A]$ has the grade and the spinor-mode sign of $A$, and that $\mathcal{L}$ is fixed by the spinor-mode sign flip.

\emph{Step 5: graded locality (clause (iii), second part).} To all orders this is \cite{DutschFredenhagen1999} (2.25) for QED and \cite{Hollands2008} (cited from memory) with coupled ghosts. At first order it follows from the free statistics, and we give the step since it shows where evenness of $\mathcal{L}$ enters. Write $[X, Y\} = XY - (-1)^{|X||Y|}YX$ for homogeneous $X$, $Y$. With the first-order retarded product $R(\mathcal{L}(y); A(x)) = \theta(x^{0} - y^{0})[A(x), \mathcal{L}(y)]$ (up to a constant), for $x$ and $z$ spacelike separated the first-order part of $[A_{\mathrm{int}}(x), B_{\mathrm{int}}(z)\}$ is a multiple of $\int g(y)\bigl(\theta(x^{0} - y^{0})[[A(x), \mathcal{L}(y)], B(z)\} + \theta(z^{0} - y^{0})[A(x), [B(z), \mathcal{L}(y)]\}\bigr)dy$. For $\mathcal{L}$ even, $\mathrm{ad}_{\mathcal{L}}$ is an even derivation of the graded commutator: $[\mathcal{L}, [A, B\}] = [[\mathcal{L}, A], B\} + [A, [\mathcal{L}, B]\}$ (graded Jacobi; checked for all pairs of five homogeneous elements in $\Cl(6, 6)$; with an odd $\mathcal{L}$ the identity needs the sign $(-1)^{|\mathcal{L}||A|}$ on the second term and the even form fails, both checked). Where both step functions are $1$ the integrand is $-[\mathcal{L}(y), [A(x), B(z)\}] = 0$, since free fields at spacelike separation graded-commute. Where only the first is $1$, $y^{0} > z^{0}$ and $[A(x), \mathcal{L}(y)] \neq 0$ only for $y$ in the causal past $J^{-}(x)$; then $y$ is spacelike to $z$ (the light-cone lemma: $y \in J^{+}(z)$ would give $x \in J^{+}(z)$, and $y \in J^{-}(z)$ contradicts $y^{0} > z^{0}$; the script confirms it on a rational grid of trial points for three spacelike pairs, and finds grid points in $J^{-}(x) \cap J^{-}(z)$ with $y^{0} \le z^{0}$, so the hypothesis $y^{0} > z^{0}$ is needed), so $B(z)$ graded-commutes with $A(x)$ and with $\mathcal{L}(y)$ and hence, by the Jacobi identity, with $[A(x), \mathcal{L}(y)]$. The last case is symmetric. So the first-order bracket vanishes. Together with Step 4, on interacting fields of even ghost count the statistics grading is $U(2\pi)$ to every order; a ghost field of odd ghost count has grading $-U(2\pi)$, the Krein case of the free-field result.

\emph{Step 6: scope (clause (iv)).} The statement is about formal power series in the coupling; the interacting local algebras exist without the adiabatic limit \cite{DutschFredenhagen1999}, but convergence and a nonperturbative vacuum are not addressed. The free-field input is the store's result for massless fields; the series' masses come from Yukawa couplings to the rung field and its vacuum value, which are even two-spinor insertions covered by Step 2, so the expansion can be posed around massless free fields (massive free Dirac locality is standard, cited from memory). Gauge anomalies and the positivity of the physical quotient (\cite{DutschFredenhagen1999}, Theorem 4) are not used: statistics and covariance do not depend on the Ward identities. Beyond perturbation theory the spin-statistics connection rests on a premise: the Jost-point theorem of Wightman's framework (Fact \claimref{field-theory.imported-standard-field-theory}) or a modular premise (\claimref{synthesis.wedge-reflections-compose-to-the-two-pi-rotor-so-geometric-modular-conjugations-give-interacting-spin-statistics}). For electrically charged fields the nonperturbative charge classes are localised in hypercones and are of Bose or Fermi type with a covariant positive-energy representation \cite{BuchholzRoberts2014} (abstract and introduction read); a spin-statistics theorem for them was not found there. Transport-absent: the inheritance is standard causal perturbation theory; the series supplies the field list, the spacetime-algebra form of its monomials and the tower ghosts.
